% TOP_PROBLEM: {"problemId": "problem.top500-omission-w045049-bernoulli-absolute-continuity-non-pisot", "canonicalTitle": "Absolute continuity of Bernoulli convolutions outside reciprocal Pisot parameters", "releaseRank": 426, "primaryDomain": "probability_ergodic_dynamics", "primaryDomainLabel": "Probability, ergodic theory & dynamics", "displayStatus": "open_with_solved_subcases", "releaseStatus": "open_with_solved_subcases"}
% !TeX program = lualatex
% Definition and sources only; this file is not an LLM solution attempt.
\documentclass[11pt]{article}
\usepackage[margin=25mm]{geometry}
\usepackage{fontspec}
\setmainfont{Latin Modern Roman}
\usepackage{amsmath,amssymb,mathtools}
\usepackage{unicode-math}
\setmathfont{Latin Modern Math}
\usepackage{xurl,hyperref}
\hypersetup{colorlinks=true,urlcolor=blue}
\setcounter{secnumdepth}{0}
\setlength{\parindent}{0pt}
\setlength{\parskip}{5pt}
\setlength{\emergencystretch}{3em}
\begin{document}
\section*{\texorpdfstring{Absolute continuity of Bernoulli convolutions outside reciprocal Pisot parameters}{Absolute continuity of Bernoulli convolutions outside reciprocal Pisot parameters}}\label{426}
\subsection{Definitions and mathematical statement}
For $1/2<\lambda<1$, let $(\epsilon_n)_{n\ge0}$ be independent random signs, each equally likely to be $1$ or $-1$, and let $\nu_\lambda$ be the probability distribution of the almost surely convergent series
\[
 X_\lambda=\sum_{n=0}^{\infty}\epsilon_n\lambda^n.
\]
A Pisot number is a real algebraic integer $\beta>1$ all of whose other algebraic conjugates have modulus less than one. The selected assertion is
\[
 \lambda^{-1}\text{ is not Pisot}\quad\Longrightarrow\quad
 \nu_\lambda\ll\mathcal L^1.
\]
Absolute continuity means that every Lebesgue-null Borel set has $\nu_\lambda$-measure zero, equivalently that $\nu_\lambda$ has an integrable density. The quantifier is over every eligible parameter, not almost every parameter. Full Hausdorff dimension alone would not establish absolute continuity. Multiplying $X_\lambda$ by a nonzero normalization constant leaves this question unchanged.
\par\smallskip\textit{Formulation and concept references:} \href{https://londmathsoc.onlinelibrary.wiley.com/doi/full/10.1112/jlms.70515}{[S1]}.

\subsection{Short English statement}
Must a random signed geometric series have a probability density whenever the reciprocal of its contraction parameter is not a Pisot number?
\subsection{Sources}
\begin{itemize}
\item[S1]Baker, Koivusalo, Troscheit and Zhang, On the Fourier transform of random Bernoulli convolutions, JLMS 113 (2026), e70515.
\url{https://londmathsoc.onlinelibrary.wiley.com/doi/full/10.1112/jlms.70515}.
\textit{Formulation source.}
\end{itemize}
\end{document}
